% This is part of Mes notes de mathématique
% Copyright (c) 2008-2020
%   Laurent Claessens
% See the file fdl-1.3.txt for copying conditions.

%+++++++++++++++++++++++++++++++++++++++++++++++++++++++++++++++++++++++++++++++++++++++++++++++++++++++++++++++++++++++++++
\section{Parties libres, génératrices, bases et dimension}
%+++++++++++++++++++++++++++++++++++++++++++++++++++++++++++++++++++++++++++++++++++++++++++++++++++++++++++++++++++++++++++

Nous avons déjà défini (dans~\ref{DEFooKHWZooIfxdNc}) un espace vectoriel comme étant un module sur un corps commutatif. En explicitant un peu, cela donne ceci\cite{ooQLVLooEUrNLS}.

Un espace vectoriel sur le corps \( \eK\) est un ensemble \( V\) muni de deux opérations :
\begin{itemize}
    \item une loi de composition interne \( +\colon V\times V\to V\),
    \item une loi de composition externe \( \cdot\colon \eK\times V\to V\)
\end{itemize}
telles que
\begin{enumerate}
    \item
        \( (V,+)\) soit un groupe abélien,
    \item
        pour tout \( u,v\in V\) et pour tout \( k,k'\in \eK\),
        \begin{subequations}
           \begin{align}
                k(u+v)=(ku)+(kv)\\
                (kk')u=k(k'u)\\
                (k+k')u=(ku)+(k'u)\\
                1u=u
            \end{align}
        \end{subequations}
        où \( 1\) est le neutre de \( \eK\) et où nous avons directement adopté la notation \( ku\) pour \( k\cdot u\).
\end{enumerate}
Si \( u\in V\), nous notons \( -u\) l'inverse de \( u\) dans le groupe \( (V,+)\).

\begin{definition}[Partie libre]
    Si \( E\) est un espace vectoriel, une partie \( A\) de \( E\) est \defe{libre}{libre!partie} si pour tout choix d'un nombre fini d'éléments \( \{ u_i \}_{i=1,\ldots, n}\), l'égalité
    \begin{equation}
        a_1 u_1+\cdots +a_nu_n=0
    \end{equation}
    implique \( a_i=0\) pour tout \( i\) (ici les \( a_i\) sont dans le corps de base).

    Une partie infinie est libre si toutes ses parties finies le sont.
\end{definition}

\begin{remark}
    Notons que le vecteur nul n'est dans aucune partie libre, ne fût-ce que parce que \( a0=0\) n'implique pas \( a=0\).
\end{remark}

Si \( A\) est une partie de l'espace vectoriel \( E\) nous notons \( \Span(A)\)\nomenclature[A]{$\Span(A)$}{l'ensemble des combinaisons linéaires finies d'éléments de \( A\)} l'ensemble des combinaisons linéaires finies d'éléments de \( A\). Les coefficients de ces combinaisons linéaires sont dans le corps de base \( \eK\).

\begin{definition}[Partie génératrice]
    Une partie $B$ d'un espace vectoriel \( E\) est \defe{génératrice}{partie!génératrice} si \( \Span(B)=E\).
\end{definition}

\begin{remark}
    Ces définitions demandent des commentaires en dimension infinie\footnote{Nous n'avons pas encore défini le concept de dimension, mais nous nous adressons au lecteur trop pressé.}.

    \begin{enumerate}
        \item
    Tout élément peut être écrit comme combinaison linéaire finie d'une partie génératrice. Cela ne signifie pas que nous pouvons extraire une partie finie qui convient pour tous les éléments à la fois. Lorsque l'espace est de dimension infinie, ceci est particulièrement important.
\item
    La définition séparée de liberté dans le cas des parties infinies a son importance lorsqu'on parle d'espaces vectoriels de dimension infinies (en dimension finie, aucune partie infinie n'est libre) parce que cela fera une différence entre une base algébrique et une base hilbertienne par exemple.
    \end{enumerate}
\end{remark}

\begin{definition}[Base]        \label{DEFooNGDSooEDAwTh}
    Une \defe{base}{base} de l'espace vectoriel \( E\) est une partie à la fois génératrice et libre.
\end{definition}

\begin{proposition}[\cite{MonCerveau}]      \label{PROPooEIQIooXfWDDV}
    Tout élément non nul d'un espace vectoriel possédant une base\footnote{Nous n'avons pas démontré que tout espace vectoriel possède une base. Donc à notre niveau, il est possible que ce théorème soit sans objet pour beaucoup d'espaces.} se décompose de façon unique en combinaison linéaire finie d'éléments d'une base.
\end{proposition}

\begin{proof}
    Soit un espace vectoriel \( E\) et une base \( \{ e_i \}_{i\in I}\) où \( I\) est un ensemble a priori quelconque. Soit \( v\in E\). Vu que \( E=\Span\{ e_i \}_{i\in I}\), il existe une partie finie \( J\) de \( I\) et des coefficients \( \{ v_j \}_{j\in J}\) dans \( \eK\) tels que
    \begin{equation}
        v=\sum_{j\in J}v_je_j.
    \end{equation}
    Cela donne l'existence.

    En ce qui concerne l'unicité, soient \( J \) et \( K\) des parties finies de \( I\) et des coefficients \( \{ v_j \}_{j\in J}\) et \( \{ w_{k} \}_{k\in K}\) tels que
    \begin{equation}
        v=\sum_{j\in J}v_je_j=\sum_{k\in K}w_{k}e_{k}.
    \end{equation}
    Nous posons \( L=J\cup K\) et, pour \( l\in L\),
    \begin{equation}
        \alpha_l=\begin{cases}
            v_l    &   \text{si } l\in J\setminus K\\
            w_l    &    \text{si } l\in K\setminus J\\
            v_l-w_l    &    \text{si } l\in K\cap J.
        \end{cases}
    \end{equation}
    Nous avons alors
    \begin{equation}
        \sum_{l\in L}\alpha_le_l=0,
    \end{equation}
    ce qui implique que \( \alpha_l=0\) pour tout \( l\in L\) parce que la partie \( \{ e_i \}_{i\in I}\) est libre et que \( L\) est finie.

    L'unicité de la décomposition de \( v\) signifie que
    \begin{equation}
        \{ j\in J \tq v_j\neq 0 \}=\{ k\in K\tq w_k\neq 0 \}
    \end{equation}
    et que pour \( l\) dans cet ensemble, \( v_l=w_l\).

    Soit \( j\in J\); il y a deux possibilités : \( j\in J\setminus K\) et \( j\in J\cap K\). Dans le premier cas nous avons déjà vu que \( \alpha_j=v_j=0\). Dans le second cas, \( \alpha_j=v_j-w_j=0\), c'est-à-dire \( v_j=w_j\).

    Donc \( j\in J\) vérifiant \( v_j\neq 0\) implique \( j\in J\cap K\) et l'égalité des coefficients. Idem avec \( k\in K\) tel que \( w_k\neq 0\) implique \( k\in J\cap K\).
\end{proof}

\begin{definition}
    Un espace vectoriel est \defe{de type fini}{type!fini!espace vectoriel} s'il contient une partie génératrice finie.
\end{definition}
Nous verrons dans les résultats qui suivent que cette définition est en réalité inutile parce qu'une espace vectoriel sera de type fini si et seulement s'il est de dimension finie.

\begin{lemma}       \label{LemytHnlD}
    Si \( E\) a une famille génératrice de cardinal \( n\), alors toute famille de \( n+1\) éléments est liée.
\end{lemma}

\begin{proof}
    Nous procédons par récurrence sur \( n\). Pour \( n=1\), nous avons \( E=\Span(e)\) et donc si \( v_1,v_2\in E\) nous avons \( v_1=\lambda_1 e\), \( v_2=\lambda_2e\) pour certains éléments non nuls \( \lambda_1,\lambda_2\) du corps de base. Nous avons donc \( \lambda_2v_1-\lambda_1v_1=0\). Cela prouve que \( \{ v_1,v_2 \}\) est liée.

    Supposons maintenant que le résultat soit vrai pour \( k<n\), c'est-à-dire que pour tout espace vectoriel contenant une partie génératrice de cardinal \( k<n\), les parties de \( k+1\) éléments sont liées. Soit maintenant un espace vectoriel muni d'une partie génératrice \( G=\{ e_1,\ldots, e_n \}\) de \( n\) éléments, et montrons que toute partie \( V=\{ v_1,\ldots, v_{n+1} \}\) contenant \( n+1\) éléments est liée. On peut supposer que:
    \begin{itemize}
      \item
        les \( v_i \) sont tous non-nuls (sinon, un \( v_i \) nul, cela entraîne pour un \( \alpha \in \eK \) non nul, \( \alpha v_i = 0 \));
      \item
        les \( v_i \) sont tous distincts (sinon, \( v_i - v_j  = 0 \) pour ceux qui seraient égaux).
    \end{itemize}
    Étant donné que \( \{ e_i \}\) est génératrice nous pouvons définir les nombres \( \lambda_{ij}\) par
    \begin{equation}
        v_i=\sum_{k=1}^n\lambda_{ij}e_j
    \end{equation}
    Vu que
    \begin{equation}
        v_{n+1}=\sum_{k=1}^n\lambda_{n+1,k}e_k\neq 0,
    \end{equation}
    quitte à changer la numérotation des \( e_i\) nous pouvons supposer que \( \lambda_{n+1,n}\neq 0\). Nous considérons les vecteurs
    \begin{equation}
        w_i=\lambda_{n+1,n}v_i-\lambda_{i,n}v_{n+1}.
    \end{equation}
    En calculant un peu,
    \begin{subequations}
        \begin{align}
            w_i&=\lambda_{n+1,n}\sum_k\lambda_{i,k}e_k-\lambda_{i,n}\sum_k\lambda_{n+1,k}e_k\\
            &=\sum_{k=1}^{n-1}\big( \lambda_{n+1,n}\lambda_{i,k}-\lambda_{i,n}\lambda_{n+1,} \big)e_k
        \end{align}
    \end{subequations}
    parce que les termes en \( e_n\) se sont simplifiés. Donc la famille \( \{ w_1,\ldots, w_n \}\) est une famille de \( n\) vecteurs dans l'espace vectoriel \( \Span\{ e_1,\ldots, e_{n-1} \}\); elle est donc liée par l'hypothèse de récurrence. Il existe donc des nombres \( \alpha_1,\ldots, \alpha_n\in \eK\) non tous nuls tels que
    \begin{equation}        \label{EqOQGGoU}
        0=\sum_{i=1}^n\alpha_iw_i=\sum_{i=1}^n\alpha_i\lambda_{n+1,n}v_i-\left( \sum_{i=1}^n\alpha_i\lambda_{i,n} \right)v_{n+1}.
    \end{equation}
    Vu que \( \lambda_{n+1,n}\neq 0\) et que parmi les \( \alpha_i\) au moins un est non nul, nous avons au moins un des produits \( \alpha_i\lambda_{n+1,n}\) qui est non nul. Par conséquent \eqref{EqOQGGoU} est une combinaison linéaire nulle non triviale des vecteurs de \( \{ v_1,\ldots, v_{n+1} \}\). Cette partie est donc liée.
\end{proof}

\begin{lemma}   \label{LemkUfzHl}
    Soit \( L\) une partie libre et \( G\) une partie génératrice. Si l'ensemble des parties libres \( L'\) telles que \( L\subset L'\subset G\) possède un élément maximum \( B \)\footnote{Encore une fois, à part quelques cas triviaux, il n'est pas clair à ce point que ce maximum existe.}, alors cet élément est une base.
\end{lemma}
Qu'entend-on par «maximale» ? La partie \( B\) doit être libre, contenir \( L\), être contenue dans \( G\) et de plus avoir la propriété que \( \forall x\in G\setminus B\), la partie \( B\cup\{ x \}\) est liée.

\begin{proof}
    D'abord si \( G\) est une base, alors toutes les parties de \( G\) sont libres et le maximum est \( B=G\). Dans ce cas le résultat est évident. Nous supposons donc que \( G\) est liée.

    La partie \( B \) est libre parce qu'on l'a prise parmi les libres. Montrons que \( B\) est génératrice. Soit \( x\in G\setminus B\); par hypothèse de maximalité, \( B\cup\{ x \}\) est liée, c'est-à-dire qu'il existe des éléments \( b_i \in B \) et des nombres \( \lambda_i\), \( \lambda_x\) non tous nuls, tels que
    \begin{equation}    \label{EqxfkevM}
        \sum_{i=1}^l\lambda_ib_i+\lambda_xx=0.
    \end{equation}
    Si \( \lambda_x=0\) alors un des \( \lambda_i\) doit être non nul et l'équation \eqref{EqxfkevM} devient une combinaison linéaire nulle non triviale des \( b_i\), ce qui est impossible parce que \( B\) est libre. Donc \( \lambda_x\neq 0\) et
    \begin{equation}
        x=\frac{1}{ \lambda_x }\sum_{i=1}^l\lambda_ib_i.
    \end{equation}
    Donc tous les éléments de \( G\setminus B\) sont des combinaisons linéaires des éléments de \( B\), et par conséquent, \( G\) étant génératrice, tous les éléments de \( E\) sont combinaisons linéaires d'éléments de \( B\).
\end{proof}

\begin{theorem}[Théorème de la base incomplète] \label{ThonmnWKs}
    Soit \( E\) un espace vectoriel de type fini sur le corps \( \eK\).
    \begin{enumerate}
        \item     \label{ItemBazxTZ}
            Si \( L\) est une partie libre et si \( G\) est une partie génératrice contenant \( L\), alors il existe une base \( B\) telle que \( L\subset B\subset G\).
        \item     \label{ITEMooFVJXooGzzpOu}
            Toute partie libre peut être étendue en une base.
        \item     \label{ITEMooFBUAooSSZxgx}
            Toutes les bases sont finies et ont même cardinal.
        \item       \label{ITEMooJIJSooGuJMdt}
            Si \( V\) est un sous-espace vectoriel de \( E\), et si \( L\) est une base de \( V\), alors il existe une base \( E\) qui contient \( L\).
    \end{enumerate}
\end{theorem}
\index{théorème!base incomplète}

\begin{proof}
    Point par point.
    \begin{enumerate}
        \item
    Vu que \( E\) est de type fini, il admet une partie génératrice \( G\) de cardinal fini \( n\). Donc une partie libre est de cardinal au plus \( n\) par le lemme~\ref{LemytHnlD}. Soit \( L\), une partie libre contenue dans \( G\) (ça existe : par exemple \( L=\emptyset\)). La partie \( B\) maximalement libre contenue dans \( G\) et contenant \( L\) est une base par le lemme~\ref{LemkUfzHl}.
\item
Notons que puisque \( E\) lui-même est générateur, le point~\ref{ItemBazxTZ} implique que toute partie libre peut être étendue en une base.
\item
    Soient \( B\) et \( B'\), deux bases. En particulier \( B\) est génératrice et \( B'\) est libre, donc le lemme~\ref{LemytHnlD} indique que \( \Card(B')\leq \Card(B)\). Par symétrie on a l'inégalité inverse. Donc \( \Card(B)=\Card(B')\).
\item
    La partie \( L\) étant une base de \( V\), elle est en particulier libre dans \( E\). Par le point \ref{ITEMooFVJXooGzzpOu}, \( L\) peut être étendue en une base.
    \end{enumerate}
\end{proof}

\begin{remark}      \label{REMooYGJEooEcZQKa}
    Le théorème de la base incomplète~\ref{ThonmnWKs}\ref{ITEMooFVJXooGzzpOu} est ce qui permet de construire une base d'un espace vectoriel en « commençant par» une base d'un sous-espace. En effet si \( H\) est un sous-espace de \( E\) alors une base de \( H\) est une partie libre de \( E\) et donc peut être étendue en une base de \( E\).
\end{remark}

\begin{definition}      \label{DEFooWRLKooArTpgh}
    La \defe{dimension}{dimension} d'un espace vectoriel \( E \) de type fini, noté \(\dim E \), est le cardinal\footnote{Définition \ref{PROPooJLGKooDCcnWi}.} d'une\footnote{Le théorème de la base incomplète~\ref{ThonmnWKs}\ref{ITEMooFBUAooSSZxgx} montre que cette définition ne souffre d'aucune ambiguïté.} de ses bases.
\end{definition}
\index{dimension!définition}

Il existe une infinité de bases de $\eR^m$. On peut démontrer que le cardinal de toute base de $\eR^m$ est $m$, c'est-à-dire que toute base de $\eR^m$ possède exactement $m$ éléments.

\begin{example}
    Sur l'espace vectoriel \( \eR^m\) sur \( \eR \), on considère pour tout $j = 1,\dots, m$, le vecteur $e_j$ défini par
    \begin{equation}\nonumber
      e_j=
    \begin{array}{cc}
      \begin{pmatrix}
        0\\\vdots\\0\\1\\ 0\\\vdots\\0
      \end{pmatrix} &
      \begin{matrix}
        \quad\\\quad\\\leftarrow\textrm{j-ème} \quad\\\quad\\\quad\\
      \end{matrix}
    \end{array}.
    \end{equation}
    La partie $\{e_1,\ldots, e_m\}$ est alors une base, appelée \defe{base canonique}{canonique!base}\index{base canonique de $\eR^m$} de \( \eR^m\).
    La composante numéro $j$ de $e_i$ est $1$ si $i=j$ et $0$ si $i\neq j$. Cela s'écrit $(e_i)_j=\delta_{ij}$ où $\delta$ est le \defe{symbole de Kronecker}{Kronecker} défini par
    \begin{equation}
        \delta_{ij}=\begin{cases}
            1	&	\text{si }i=j\\
            0	&	 \text{si }i\neq j
        \end{cases}
    \end{equation}
    Les éléments de la base canonique de $\eR^m$ peuvent donc être écrits $e_i=\sum_{k=1}^m\delta_{ik}e_k$.
\end{example}

Le théorème suivant est essentiellement une reformulation du théorème~\ref{ThonmnWKs}.
\begin{theorem} \label{ThoMGQZooIgrXjy}
    Soit \( E\) un espace vectoriel de dimension finie et \( \{ e_i \}_{i\in I}\) une partie génératrice de \( E\).

    \begin{enumerate}
        \item       \label{ITEMooTZUDooFEgymQ}
            Il existe \( J\subset I\) tel que \( \{ e_i \}_{i\in J}\) est une base. Autrement dit : de toute partie génératrice nous pouvons extraire une base.
        \item       \label{ITEMooCJQGooXwjsfm}
            Soit \( \{ f_1,\ldots, f_l \}\) une partie libre. Alors nous pouvons la compléter en utilisant des éléments \( e_i\). C'est-à-dire qu'il existe \( J\subset I\) tel que \( \{ f_k \}\cup\{ e_i \}_{i\in J}\) soit une base.
    \end{enumerate}
\end{theorem}

\begin{proposition}     \label{PROPooVEVCooHkrldw}
    Si \( E\) est un espace vectoriel de dimension finie \( n\), alors 
    \begin{enumerate}
        \item       \label{ITEMooZNLDooBISkJyBS}
            toute partie contenant \( n+1\) éléments est liée.
        \item       \label{ITEMooSGGCooOUsuBs}
            toute partie libre contenant \( n\) éléments est une base,
        \item
            toute partie génératrice contenant \( n\) éléments est une base.
    \end{enumerate}
\end{proposition}

\begin{proof}
    Soit une partie \( M\) contenant \( n+1\) éléments. L'espace \( E\) possède une partie génératrice contenant \( n\) éléments (n'importe quelle base). Donc \( M\) est liée par le lemme \ref{LemytHnlD}.

    Une partie libre contenant \( n\) éléments peut être étendue en une base; si ladite extension est non triviale (c'est-à-dire qu'on ajoute vraiment au moins un élément) une telle base contiendra une partie de \( n+1\) éléments qui serait liée par le lemme~\ref{LemytHnlD}.

    Pour l'autre assertion, soit une partie génératrice \( \{ v_i \}_{i\in I}\) où \( I\) contient \( n\) éléments. Par le théorème \ref{ThoMGQZooIgrXjy}\ref{ITEMooCJQGooXwjsfm} il existe \( J\subset I\) tel que \( \{ v_j \}_{j\in J}\) soit une base. Si l'inclusion \( J\subset I\) est stricte, alors la base \( \{ v_j \}_{j\in J}\) contiendrait moins de \( n\) éléments, ce qui serait en contradiction avec le théorème \ref{ThonmnWKs}\ref{ITEMooFBUAooSSZxgx}.
\end{proof}

\begin{proposition}[\cite{ooDSTAooKgSyCN}]      \label{PROPooQCIXooHIyPPq}
    Soit \( E\), un espace vectoriel de dimension finie sur le corps $\eK$. Soient \( V\) et \( W\) des sous-espaces vectoriels de \( E\). Alors
    \begin{equation}
        \dim(V+W)=\dim(V)+\dim(W)-\dim(V\cap W).
    \end{equation}
\end{proposition}

\begin{proof}
    Commençons par remarquer que, comme \( E \) est de dimension finie, tous les sous-espaces vectoriels sont aussi de dimension finie. Considérons le sous-espace vectoriel \( V \cap W \). Deux possibilités: soit il est réduit à \( \{0\}\), soit il a au moins un élément non nul. Dans les deux cas, on pose \( m = \dim(V \cap W) \), et on fixe \( \{e_i\}_{i=1}^m \) une base de \( V \cap W\), éventuellement vide. (Quand elle n'est pas vide, elle existe, grâce au théorème de la base incomplète~\ref{ThonmnWKs}~\ref{ItemBazxTZ}). La famille  \( \{e_i\}_{i=1}^m \) est libre\footnote{Qu'en est-il si la famille est vide? Pure question de logique: elle est libre.} dans \( V \) et dans \( W \); on l'étend dans chacun de ces espaces par respectivement des \( v_j \) (on en a \( \dim(V) - m \)) et des \( w_k \) (on en a \( \dim(W) - m \)). Le nombre d'éléments ajoutés est connu, par le théorème~\ref{ThonmnWKs}~\ref{ITEMooJIJSooGuJMdt}, et la définition de dimension~\ref{DEFooWRLKooArTpgh}
    
    Nous allons démontrer que la famille \( B \) constituée des \( e_i \), des \( v_j \) et des \(w_k \), forme une base de \( V + W \). Dès lors, on aura l'égalité suivante:
    \begin{align}
      \dim(V+W) &= m + (\dim(V) - m) + (\dim(W) - m )\\
      & = \dim(V)+\dim(W) - m = \dim(V)+\dim(W)-\dim(V\cap W).
    \end{align}
    La famille \(B \) est clairement génératrice de \( V + W \): en effet, si \( v + w \in V + W \), alors on décompose \( v \in V \) dans la base des  \( e_i \) et des \( v_j \), et \( w \in W \) dans la base des \( e_i \) et des \( w_k \), puis on additionne. Reste à voir qu'elle est libre.
    
    On considère \( \alpha_i \in \eK \) tels que \( \sum_i \alpha_i e_i + \sum_j \alpha_j v_j + \sum_k \alpha_k w_k = 0 \). Alors,
    \begin{equation}
         \sum_i \alpha_i e_i + \sum_j \alpha_j v_j = - \sum_k \alpha_k w_k.
    \end{equation}
    Le membre de gauche est dans \( V \), celui de droite est dans \( W \). Donc nécessairement, cette égalité a lieu dans \( V \cap W\), un endroit où ni les \( v_j \), ni les \( w_k \) n'ont droit de cité: les coefficients \( \alpha_j \) et \( \alpha_k \) qui apparaissent ci-dessus sont nuls. Par suite,
        \begin{equation}
         \sum_i \alpha_i e_i = 0
    \end{equation}
    dans \( V \cap W \), donc tous les \( \alpha_i \) aussi sont nuls, car les \( e_i \) forment une base de \( V \cap W \).
\end{proof}


\begin{definition}\label{DefCodimension}
Soit \( F\) un sous-espace vectoriel de l'espace vectoriel \( E\). La \defe{codimension}{codimension} de \( F\) dans \( E\) est
\begin{equation}
    \codim_E(F)=\dim(E/F).
\end{equation}
\end{definition}

\begin{probleme}
Voir que $E/F$ a une structure vectorielle, expliciter sa dimension en fonction de celles de $E$ et $F$. A priori \( \dim E - \dim F \), mais il faut le justifier proprement.
\end{probleme}
%---------------------------------------------------------------------------------------------------------------------------
\subsection{Et en dimension infinie}
%---------------------------------------------------------------------------------------------------------------------------

Dans ZFC, en dimension infinie, il existe aussi une base pour tout espace vectoriel ainsi qu'un théorème de la base incomplète. Nous ne parlerons pas de ce qu'il se passe lorsque nous ne considérons que ZF\footnote{Si vous ne savez pas ce que signifient les sigles «ZF» et «ZFC» vous ne devriez pas être en train de lire ceci, et encore moins penser à le resservir à un jury d'agrégation.}.

\begin{lemma}[\cite{ooXEFKooHikcdE}]        \label{LEMooSSRXooIyfgNz}
    Soient un \( \eK\)-espace vectoriel \( E\) et un sous-espace vectoriel \( V\) de \( E\). Soient encore deux sous-espaces vectoriels \( W_1\) et \( W_2\) tels que
    \begin{enumerate}
        \item
            \( V\cap W_1=\{ 0 \}\);
        \item
            \( V+W_2=E\).
    \end{enumerate}
    Alors il existe un supplémentaire \( W\) de \( V\) tel que \( W_1\subset W\subset W_2\).
\end{lemma}

Juste une remarque : dans le Frido le symbole «\( \subset\)» ne signifie pas une inclusion stricte.

\begin{proof}
    Nous divisons en petits morceaux.
    \begin{subproof}
        \item[Un gros ensemble]
            Soit \( \mA\) l'ensemble des sous-espaces vectoriels \( S\) de \( E\) tels que \( W_1\subset S\subset W_2\) et \( S\cap V=\{ 0 \}\). Vu que \( W_1\subset \mA\), cet ensemble n'est pas vide. De plus \( \mA\) est partiellement ordonné pour l'inclusion.
        \item[\( \mA\) est inductif]
            Nous prouvons maintenant que \( \mA\) est inductif\footnote{Définition~\ref{DefGHDfyyz}.}. Pour cela, soit une partie \( \mA'\) totalement ordonnée et \( U=\bigcup_{A\in \mA'}A\).

            Alors, la partie \( U\) est un sous-espace vectoriel de \( E\). En effet si \( x,y\in U\), alors il existe \( A_1,A_2\in\mA'\) tels que \( x\in A_1\) et \( y\in A_2\). Vu que \( \mA'\) est totalement ordonné, l'un des ensembles parmi \( A_1\) et \( A_2\) est inclus dans l'autre. Sans perdre de généralité, disons \( A_1\subset A_2\). Alors les opérations s'effectuent dans \( A_2 \) : nous avons \( x,y\in A_2\), et donc \( x+y\in A_2\subset U\). La multiplication par un scalaire est plus simple:  \( \lambda x\in A_1\subset U\).

            La sous-espace vectoriel \( U \) contient \( W_1 \), il est contenu dans \( W_2\), et il ne contient pas d'élément de \( V \) autre que \( 0 \). Ainsi, \( U\in \mA\) et majore \( \mA'\) pour l'inclusion. En bref, \( \mA\) est bien inductif.
        \item[Utilisation de Zorn]

            Le lemme de Zorn~\ref{LemUEGjJBc} nous donne alors un maximum \( W\) de \( \mA\). Ce maximum vérifie
            \begin{enumerate}
                \item
                    \( W\cap V=\{ 0 \}\),
                \item
                    \( W_1\subset W\subset W_2\),
                \item
                    pour tout \( W'\in\mA\), nous avons \( W'\subset W\) parce que \( W\) est maximum.
            \end{enumerate}
        \item[Supplémentaire]
            Montrons que ce \( W\) est un supplémentaire de \( V\). Soit \( x\in E\). Le but est de trouver une décomposition de \( x\) en somme d'un élément de \( W\) et un de \( V\). Vu que \( V+W_2=E\) nous avons \( v\in V\) et \( w_2\in W_2\) tels que 
            \begin{equation}
                x=v+w_2. 
            \end{equation}
            Si \( w_2\in W\) alors c'est fait. Sinon \ldots

            Soit \( X=\Span\{ W,w_2 \}\). Vu que \( X\) contient strictement \( W\) et que \( W\) est maximum dans \( \mA\), la partie \( X\) n'est pas un élément de \( \mA\). Vu que \( X\) est un sous-espace vectoriel de \( E\) tel que \( W_1\subset X\subset W_2\), la seule possibilité pour que \( X\) ne soit pas dans \( \mA\) est que \( X\cap V\neq \{ 0 \}\). Soit donc \( y\neq 0\) dans \( X\cap V\). Par définition de \( X\),
            \begin{equation}\label{EqDecompo55:296}
                y=w'+\lambda w_2
            \end{equation}
            pour \( w'\in W\), \( w_2\in W_2\) et \( \lambda\in \eK\). Nous avons \( \lambda\neq 0\), sinon nous aurions \( y\in W\cap V \) et donc \(y = 0 \) puisque \( W \) est dans \( \mA \). La décomposition \eqref{EqDecompo55:296} permet alors d'écrire \( w_2=(y-w')/\lambda\) et finalement
            \begin{equation}
                x=v+\frac{1}{ \lambda }(y-w')=\underbrace{v+\frac{1}{ \lambda }y}_{\in V}-\underbrace{\frac{1}{ \lambda }w'}_{\in W}.
            \end{equation}
            La somme d'espaces vectoriels \( E=V+W\) est donc établie.
    \end{subproof}
\end{proof}

\begin{corollary}
    Tout sous-espace vectoriel d'un espace vectoriel possède un supplémentaire.
\end{corollary}

\begin{proof}
    Soit un espace vectoriel \( E\) ainsi qu'un sous-espace vectoriel \( V\). Si \( V=E\) nous sommes ok. Sinon nous considérons \( v\in E\setminus V\) et nous posons \( W_1=\eK v\) et \( W_2=E\).

    Vu que \( V\) et \( W_1\) sont des espaces vectoriels, nous avons \( V\cap W_1=\{ 0 \}\), et vu que \( W_2=E\) nous avons \( V+W_2=E\). Le lemme~\ref{LEMooSSRXooIyfgNz} nous donne alors un supplémentaire de \( V\).
\end{proof}

\begin{proposition}[Base incomplète]        \label{PROPooHDCEooMhDjPi}
    Tout espace vectoriel (non réduit à \( \{ 0 \}\)) possède une base.
\end{proposition}

\begin{proof}
    Soit \( \mA\) l'ensemble des familles libres de \( E\). Il n'est pas vide parce que \( \{ v \}\) en est une dès que \( v\) est non nul dans \( E\). Rapidement :
    \begin{itemize}
        \item l'ensemble \( \mA\) est ordonné pour l'inclusion,
        \item si \( \mA'\) est une partie totalement ordonnée, l'union est un majorant,
        \item donc \( \mA\) est inductif,
        \item soit un maximum \( F\) de \( \mA\).
    \end{itemize}
    La partie \( F\) est libre parce qu'elle est dans \( \mA\). Elle est génératrice parce que si \( v\) n'est pas dans \( \Span(F)\) alors la partie \( F\cup\{ v \}\) est encore libre, et majore strictement $F$ pour l'inclusion, ce qui n'est pas possible.

    Donc \( F\) est une base de \( E\).
\end{proof}

\begin{theorem}[Base incomplète, dimension quelconque]      \label{THOooOQLQooHqEeDK}
    Soit une partie \( \{ e_i \}_{i\in I}\) génératrice de l'espace vectoriel \( E\) (ici, \( I\) est un ensemble quelconque\footnote{Un cas d'utilisation intéressant est de poser \( I=E\) et \( e_i=i\). Pensez-y.}). Soit \( I_0\subset I\) tel que \( \{ e_i \}_{i\in I_0}\) soit libre.

    Alors il existe \( I_1\) tel que \( I_0\subset I_1\subset I\) tel que \( \{ e_i \}_{i\in I_1}\) soit une base de \( E\).
\end{theorem}

Note : une telle partie \( I_0\) existe en prenant un singleton. Mais l'existence n'est pas le sujet ici.

\begin{proof}
    Soit \( \mA\) l'ensemble des parties \( J\) de \( I\) telles que \( I_0\subset J\subset I\) et telles que \( \{ e_i \}_{i\in J}\) soit libre.

    Encore une fois, \( \mA\) est inductif pour l'ordre partiel donné par l'inclusion. Soit \( J\) un maximum. Vu que \( J\in\mA\), la partie \( \{ e_i \}_{i\in J}\) est libre. Mais elle est également génératrice parce que si \( e_k\) n'est pas dedans, \( J\) ne serait pas maximum, étant majorée par \( J\cup\{ k \}\).

    Donc \( \{ e_i \}_{i\in J}\) engendre tous les \( e_i\) avec \( i\in I\) et donc tous les éléments de \( E\).
\end{proof}

%--------------------------------------------------------------------------------------------------------------------------- 
\subsection{Espace librement engendré}
%---------------------------------------------------------------------------------------------------------------------------

\begin{definition}[\cite{ooGNYOooGZKGba}]       \label{DEFooCPNIooNxsYMY}
    Soient un ensemble \( S\) et un corps \(\eK \). L'espace vectoriel \defe{librement engendré}{librement engendré} sur \( S\), noté \( F_{\eK}(S)\) est l'ensemble des applications \( S\to \eK\) qui sont non-nulles en un nombre fini de points de \( S\).

    Autrement dit, \( \sigma\colon S\to \eK\) est dans \( F_{\eK}(S) \) si \( \{ x\in S\tq \delta(x)\neq 0 \}\) est fini\footnote{Parce que nous l'aimons bien, nous ne résistons pas à faire un renvoi vers la définition \ref{DefEOZLooUMCzZR}.}.
\end{definition}

Le lemme suivant donne tout son sens à l'expression «librement» engendré. Il dit que \( F(S)\) possède une base indexée par \( S\) lui-même.
\begin{lemma}       \label{LEMooLOPAooUNQVku}
    L'ensemble des applications \( \delta_s\) données par
    \begin{equation}
        \begin{aligned}
            \delta_s\colon S&\to \eK \\
            t&\mapsto \begin{cases}
                1    &   \text{si } t=s\\
                0    &    \text{sinon }
            \end{cases}
        \end{aligned}
    \end{equation}
    avec \( s\in S\) forment une base\footnote{Définition \ref{DEFooNGDSooEDAwTh}.} de \( F(S)\).
\end{lemma}

\begin{proof}
    Pour prouver que les \( \delta_s\) sont générateurs, nous considérons \( g\colon S\to \eK\) non nul sur la partie finie \( \{ s_i \}_{i\in I}\) de \( S\). Alors nous avons
    \begin{equation}
        g=\sum_{i\in I}g(s_i)\delta_{s_i}.
    \end{equation}
    
    Pour prouver que les \( \delta_s\) forment une partie libre, nous supposons avoir \( \lambda_i\in \eK\) tels que
    \begin{equation}
        g=\sum_{i\in I}\lambda_i\delta_{s_i}=0
    \end{equation}
    Soit \( j\in I\). Nous avons
    \begin{equation}
        0=f(s_j)=\sum_{i\in I}\lambda_i \underbrace{\delta_{s_i}(s_j)}_{=\delta_{ij}}=\lambda_j.
    \end{equation}
    Donc les coefficients \( \lambda_i\) sont tous nuls, et nous avons prouvé que la partie est libre.
\end{proof}

Il est parfois pratique d'écrire les éléments de \( F(S)\) comme sommes «formelles» d'éléments de \( S\). Cela va encore lorsque \( S\) est un ensemble n'ayant aucune somme bien définie. 

Mais attention : si \( S=\eR\), l'élément \( 4+7\) de \( F(\eR)\) n'est pas \( 11\). L'élément \( 11\) de \( F(\eR)\) est un élément complètement différent. Bref, il n'est pas judicieux d'écrire les éléments de \( F(S)\) comme des combinaisons linéaires d'éléments de \( S\). Pour \( x\in S\) il vaut mieux écrire explicitement \( \delta_x\) que \( x\). La somme \( \delta_x+\delta_y\) est parfaitement bien définie dans l'esemble des applications de \( S\) vers \( \eK\).

%+++++++++++++++++++++++++++++++++++++++++++++++++++++++++++++++++++++++++++++++++++++++++++++++++++++++++++++++++++++++++++++
\section{Applications linéaires}
%+++++++++++++++++++++++++++++++++++++++++++++++++++++++++++++++++++++++++++++++++++++++++++++++++++++++++++++++++++++++++++++

%---------------------------------------------------------------------------------------------------------------------------
\subsection{Définition}
%---------------------------------------------------------------------------------------------------------------------------

\begin{definition}      \label{DEFooULVAooXJuRmr}
    Soient des espaces vectoriels \( E \) et \( F\) sur le corps \( \eK\). Une application \( T\colon E\to F\) est dite \defe{linéaire}{linéaire!application} si
    \begin{itemize}
        \item $T(x+y)=T(x)+T(y)$ pour tout $x$ et $y$ dans \( E\),
        \item $T(\lambda x)=\lambda T(x)$ pour tout $\lambda$ dans $\eK$ et \( x\) dans \( E\).
    \end{itemize}
\end{definition}
Si vous avez bien suivi, les égalités dans la définition~\ref{DEFooULVAooXJuRmr} sont des égalités dans \( F\).

\begin{lemmaDef} \label{DefDQRooVGbzSm}
    L'ensemble de toutes les applications linéaires de \( E\) vers \( F\) est noté \( \aL(E,F)\)\nomenclature{$\aL(E,F)$}{Ensemble des applications linéaires de $E$ dans $F$} et devient un espace vectoriel sur \( \eK\) avec les définitions suivantes :
    \begin{enumerate}
        \item
            \( (T_1+T_2)(x)=T_1(x)+T_2(x)\),
        \item
            \( (\lambda T)(x)=\lambda T(x)\).
    \end{enumerate}
\end{lemmaDef}

\begin{example}
Pour tout $b$ dans $\eR$ la fonction $T_b(x)= bx$ est une application linéaire de $\eR$ dans $\eR$. En effet,
\begin{itemize}
\item  $T_b(x+y)= b(x+y)= bx + by = T_b(x)+T_b(y)$,
\item $T_b(ax)=b(ax)= abx = a T_b(x)$.
\end{itemize}
De la même façon on peut montrer que la fonction $T_{\lambda}$ définie par $T_{\lambda}(x)=\lambda x$ est un application linéaire de $\eR^m$ dans $\eR^m$ pour tout $\lambda$ dans $\eR$ et $m$ dans $\eN$.
\end{example}

\begin{example}\label{ex_affine}
	Soit $m=n$. On fixe $\lambda$ dans $\eR$ et $v$ dans $\eR^m$. L'application $U_{\lambda}$ de $\eR^m$ dans $\eR^m$ définie par $U_{\lambda}(x)=\lambda x+v$ n'est pas une application linéaire lorsque \( v \neq 0 \), parce que si \( a \) est un réel différent de \(0 \) et \( 1 \), alors \( av \neq v \), d'où
\[
U_{\lambda}(ax)=\lambda(ax)+v\neq a(\lambda x+v) =a U_{\lambda}(x).
\]
\end{example}

\begin{example}\label{exampleT_A}
	Soit $A$ une matrice fixée de $\eM(n \times m, \eR)$. La fonction $T_A\colon \eR^m\to \eR^n$ définie par $T_A(x)=Ax$ est une application linéaire. En effet,
    \begin{itemize}
        \item  $T_A(x+y)= A(x+y)= Ax + Ay = T_A(x)+T_A(y)$,
        \item $T_A(ax)=A(ax)= a(Ax) = a T_A(x)$.
    \end{itemize}
\end{example}

On peut observer que, si on identifie $\eM(1 \times 1, \eR)$ et $\eR$, on obtient le premier exemple comme cas particulier.

\begin{definition}[Quelques ensembles d'applications linéaires]      \label{DEFooOAOGooKuJSup}
    Soient \( E\) et \( F\) des espaces vectoriels.
    \begin{itemize}
        \item
            L'ensemble des applications linéaires de \( E\) vers \( F\) est noté $\aL(E,F)$, comme déjà dit en \ref{DefDQRooVGbzSm}.
        \item Une application linéaire \( E\to E\) est un \defe{endomorphisme}{endomorphisme} de \( E\). L'ensemble des endomorphismes de \( E\) est noté \( \End(E)\)\nomenclature[B]{$\End(E)$}{les endomorphismes de \( E\)}.
        \item Un endomorphisme bijectif est un \defe{automorphisme}{automorphisme!d'espace vectoriel}. L'ensemble des automorphismes de \( E\) est noté \( \Aut(E)\)\nomenclature[B]{$\Aut(E)$}{automorphisme de l'espace vectoriel \( E\)}.
        \item
            Une application linéaire bijective \( E\to F\) est un \defe{isomorphisme}{isomorphisme!espaces vectoriels} d'espace vectoriel. L'ensemble des isomorphismes \( E\to F\) est noté\footnote{Le fait d'utiliser une notation similaire à celle des matrices inversibles n'est pas anodine: le lecteur en est sans doute conscient.} \( \GL(E,F)\).
    \end{itemize}
\end{definition}

\begin{remark}
    Les ensembles définis en~\ref{DEFooOAOGooKuJSup} concernent la structure d'espace vectoriel seulement. Lorsque nous verrons la notion d'espace vectoriel normé, nous demanderons de plus la continuité, laquelle n'est pas automatique en dimension infinie. Voir aussi les définitions~\ref{DEFooTLQUooJvknvi}.
\end{remark}

\begin{definition}
    Si \( E\) est un espace vectoriel, si \( F\) est un espace vectoriel, et si \( T\colon E \to F \) est une application, le \defe{noyau}{noyau!vers un espace vectoriel} de \( T\) est le noyau de \( T\) lorsque \( F\) est vu comme un groupe pour l'addition\footnote{Définition \ref{DEFooWBIYooGNRYOp}.}, c'est-à-dire la partie
    \begin{equation}
        \ker(T)=\{ x\in X\tq T(x)=0 \}.
    \end{equation}
\end{definition}

\begin{proposition}     \label{PROPooRLLPooKYzsJp}
    Le noyau d'une application linéaire est un sous-espace vectoriel.
\end{proposition}

\begin{proof}
    Soit une application linéaire \( T\colon E\to F\). Si \( x,y\in \ker(T)\) et si \( \lambda\in \eK\) alors
    \begin{equation}
        T(x+y)=T(x)+T(y)=0+0=0,
    \end{equation}
    donc \( x+y\in \ker(T)\), et
    \begin{equation}
        T(\lambda x)=\lambda T(x)=0,
    \end{equation}
    donc \( \lambda x\in \ker(T)\).
\end{proof}

\begin{proposition}
    Si \( E\) et \( F\) sont des espaces vectoriels de dimension finie \( n\) et si \( \{ e_i \}_{i=1,\ldots, n}\) et \( \{ f_i \}_{i=1,\ldots, n}\) sont des bases respectivement de \( E\) et \( F\), alors il existe une unique application linéaire \( T\colon E\to F\) telle que \( T(e_i)=f_i\) pour tout \( i\).
\end{proposition}

\begin{proof}
    En deux parties.\begin{subproof}
        \item[Existence]
            Soit \( v\in E\). Vu que \( \{ e_i \}\) est une base, il se décompose de façon unique en \( v=\sum_iv_ie_i\). Alors définir
            \begin{equation}
                T(v)=\sum_iv_if_i
            \end{equation}
            est une bonne définition et satisfait aux exigences.
        \item[Unicité]
            Soient \( T\) et \( U\) satisfaisant aux exigences. Alors pour tout \( i\) nous avons \( T(e_i)=U(e_i)\). Si \( v\in E\) s'écrit de la forme \( v=\sum_iv_ie_i\) alors la linéarité impose \( T(v)=\sum_iv_iT(e_i)=\sum_iv_iU(e_i)=U(v)\). Donc \( T = U\).
    \end{subproof}
\end{proof}

\begin{lemma}[\cite{MonCerveau}]       \label{LEMooJXFIooKDzRWR}
    Soient des espaces vectoriels \( V\) et \( W\) de dimension finie. Soient des bases \( \{e_i\}\) de \( V\) et \( \{f_{\alpha}\}\) de \( W\). Nous posons
    \begin{equation}
        \begin{aligned}
            \varphi_{i\alpha}\colon V&\to W \\
            v&\mapsto v_if_{\alpha} 
        \end{aligned}
    \end{equation}
    où \( v_i\) est défini par la décomposition (unique) \( v=\sum_iv_ie_i\). 

    Alors :
    \begin{enumerate}
        \item
            La partie \( \{\varphi_{i\alpha}\} \) est une base de \( \aL(V,W)\).
        \item       \label{ITEMooPMLWooNbTyJI}
            Au niveau des dimensions : \( \dim\big( \aL(V,W) \big)=\dim(V)\dim(W)\).
    \end{enumerate}
\end{lemma}

\begin{proof}
    Il faut prouver que \( \{\varphi_{i\alpha}\}\) est libre et générateur.

    \begin{subproof}
        \item[Générateur]
            Soit une application linéaire \( T\colon V\to W\). Soit \( v \in V \); en décomposant \( T(v)\) dans la base \( \{f_{\alpha}\}\), on définit des applications \( T_{\alpha} : V \to \eK \) et il vient
            \begin{equation}
                T(v)=\sum_{\alpha}T_{\alpha}(v)f_{\alpha} = \sum_{\alpha i }T_{\alpha}(v_i e_i)f_{\alpha} =\sum_{\alpha i }v_iT_{\alpha}(e_i)f_{\alpha}=\sum_{\alpha i}T_{\alpha}(e_i)\varphi_{i\alpha}(v).             
            \end{equation}
            Donc \( T\) peut être écrit comme combinaison linéaire des \( \varphi_{i\alpha}\).

        \item[Libre]
            Supposons que \( \sum_{i\alpha}a_{i\alpha}\varphi_{i\alpha}=0\) pour certains coefficients \( a_{i\alpha}\in \eK\). Nous avons, pour tout \( v\in V\) :
            \begin{equation}
                0=\sum_{i\alpha}a_{i\alpha}\varphi_{i\alpha}(v)=\sum_{i\alpha}a_{i\alpha}v_if_{\alpha},
            \end{equation}
            mais comme les \( f_{\alpha}\) forment une base, chaque terme de la somme sur \( \alpha\) est nul :
            \begin{equation}
                \sum_ia_{i\alpha}v_i=0.
            \end{equation}
            Et comme cela est valable pour tout \( v\) et donc pour tout choix de \( v_i\), nous avons \( a_{i\alpha}=0\) pour tout \( i\) et pour tout \( \alpha\).
    \end{subproof}
    La formule de dimension est simplement la cardinalité de la base trouvée; c'est la définition \ref{DEFooWRLKooArTpgh}.
\end{proof}

%--------------------------------------------------------------------------------------------------------------------------- 
\subsection{Linéarité et bases}
%---------------------------------------------------------------------------------------------------------------------------

\begin{proposition}[\cite{ooZLSSooMYdbEz}]
    Soient deux espaces vectoriels \( E\) et \( F\). Une application linéaire\footnote{Définition \ref{DEFooULVAooXJuRmr}.} \( f\colon E\to F\) est injective si et seulement si \( \ker\{ f \}=\{ 0 \}\).
\end{proposition}

\begin{probleme}
Cela est aussi valable pour les groupes: le faire.
\end{probleme}

\begin{proof}
    Nous supposons que \( f\) est injective. Si \( x\in\ker(f)\), alors \( f(x)=0\). Or \( f\) est linéaire, donc \( f(0)=0\). Nous avons donc \( f(x)=f(0)\) et donc \( x=0\) parce que \( f\) est injective.

    Dans l'autre sens, soient \( x,y\) tels que \( f(x)=f(y)\). Par linéarité de \( f\) nous avons \( f(x-y)=0\), et donc \( x-y=0\) parce que \( \ker(f)=\{0\}\). Donc \( x=y\) et \( f\) est injective.
\end{proof}

\begin{proposition}[\cite{ooZLSSooMYdbEz}]      \label{PROPooZFKZooBGLSex}
    Soit \( f\in \aL(E,F)\) où \( E\) et \( F\) sont deux espaces vectoriels.
    \begin{enumerate}
        \item   \label{ITEMooPPMEooIaZqtm}
            Si \( f\) est injective et si \( \{v_i\}_{i\in I}\) est libre, alors \( \{f(v_i)\}_{i\in I}\) est libre.
        \item   \label{ITEMooOZSPooQBrDGi}
            Si \( f\) est surjective et si \( \{v_i\}_{i\in I}\) est génératrice, alors \( \{f(v_i)\}_{i\in I}\) est génératrice.
        \item   \label{ITEMooOIEYooIfdFnv}
            Si \( f\) est une bijection, alors l'image d'une base par \( f\) est une base.
    \end{enumerate}
\end{proposition}

\begin{proof}
    En trois parties.
    \begin{subproof}
        \item[\ref{ITEMooPPMEooIaZqtm}]
            Nous devons montrer que \( \{f(v_j)\}_{j\in J}\) est libre pour tout \( J\) fini dans \( I\). Soit donc une partie finie \( J\in I\) et des scalaires\footnote{Des éléments du corps de base \( \eK\).} tels que \( \sum_{j\in J}\lambda_jf(v_j)=0\). La linéarité de \( f\) donne\footnote{Voir les propriétés de la définition \ref{DEFooULVAooXJuRmr}.}
            \begin{equation}
                f\big( \sum_{i\in J}\lambda_jv_j \big)=0.
            \end{equation}
            Par injectivité de \( f\) nous avons alors \( \sum_j\lambda_jv_j=0\). Vu que les \( v_j\) eux-même forment une partie libre, nous avons \( \lambda_j=0\) pour tout \( j\in J\).
        \item[\ref{ITEMooOZSPooQBrDGi}]
            Soit \( y\in F\). Vu que \( f\) est surjective, il existe \( x\in E\) tel que \( f(x)=y\). Étant donné que \( \{v_i\}_{i\in I}\) est générateur, il existe une partie finie \( J\subset I\) et des scalaires \( \lambda_j\in \eK\) tels que
            \begin{equation}
                x=\sum_{j\in J}\lambda_jv_j.
            \end{equation}
            En appliquant \( f\) aux deux côtés, et en tenant compte de la linéarité de \( f\),
            \begin{equation}
                y=f(x)=\sum_{j\in J}\lambda_jf(v_j),
            \end{equation}
            ce qui prouve que \( y\) est une combinaison linéaire des \( f(v_j)\).
        \item[\ref{ITEMooOIEYooIfdFnv}]
            Une base est à la fois libre et génératrice et une bijection est à la fois injective et surjective. Les deux premiers points permettent de conclure.
    \end{subproof}
\end{proof}


\begin{lemma}[\cite{MonCerveau}]        \label{LEMooDJSIooYcsvhO}
    Un endomorphisme sur un espace vectoriel est une bijection si et seulement si il change toute base en une base.
\end{lemma}

\begin{corollary}[\cite{MonCerveau}]        \label{CORooXIPKooWThOsr}
    Si \( E\) et \( F\) sont des espaces vectoriels isomorphes de dimensions finies, alors leurs dimensions sont égales.
\end{corollary}

\begin{proof}
    Vu que \( E\) et \( F\) sont isomorphes, il existe une bijection \( f\colon E\to F\). Par la proposition \ref{PROPooZFKZooBGLSex}\ref{ITEMooOIEYooIfdFnv}, l'image d'une base de \( E\) est une base de \( F\). Donc les espaces \( E\) et \( F\) ont des bases contenant le même nombre d'éléments.
\end{proof}

%---------------------------------------------------------------------------------------------------------------------------
\subsection{Rang}
%---------------------------------------------------------------------------------------------------------------------------

La proposition~\ref{DefALUAooSPcmyK} et le théorème~\ref{ThoGkkffA} sont valables également en dimension infinie; ce sera une des rares incursions en dimension infinie de ce chapitre.
\begin{propositionDef}\label{DefALUAooSPcmyK}
    L'image d'une application linéaire est un espace vectoriel. La dimension de cet espace est le \defe{rang}{rang} de ladite application linéaire.
\end{propositionDef}

\begin{proof}
    Soit une application linéaire \( f\colon E\to F\). Nous considérons \( v,w\) dans l'image de \( f\) ainsi que \( \lambda\) dans le corps de base commun à \( E\) et \( F\).

    Soient \( v_0\in E\) et \( w_0\in E\) tels que \( v=f(v_0)\) et \( w=f(w_0)\). Alors \( v+w=f(v_0+w_0)\) et \( \lambda v=f(\lambda v_0)\). Donc l'image est bien un espace vectoriel.
\end{proof}

\begin{theorem}[Théorème du rang]       \label{ThoGkkffA}
    Soient \( E\) et \( F\) deux espaces vectoriels (de dimensions finies ou non) et soit \( f\colon E\to F\) une application linéaire. 
    
   Si \( (u_s)_{s\in S}\) est une base de \( \ker(f)\) et si \( \big( f(v_t) \big)_{t\in T}\) est une base de \( \Image(f)\) alors 
   \begin{equation}
   (u_s)_{s\in s}\cup (v_t)_{t\in T}
   \end{equation}
   est une base de \( E\).
    
   En dimension finie, nous avons en plus la formule suivante :
   \begin{equation}     \label{EQooUEOQooLySRiE}
       \rang(f)+\dim\ker f=\dim E,
   \end{equation}
   c'est-à-dire que le rang\footnote{Définition~\ref{DefALUAooSPcmyK}.} de \( f\) est égal à la codimension\footnote{Définition~\ref{DefCodimension}.} du noyau.
\end{theorem}
\index{théorème!du rang}

\begin{proof}
    Nous devons montrer que
    \begin{equation}
          (u_s)_{s\in S}\cup (v_t)_{t\in T}
    \end{equation}
    est libre et générateur.

    Soit \( x\in E\). Nous définissons les nombres \( x_t\) par la décomposition de \( f(x)\) dans la base \( \big( f(v_t) \big)\) :
    \begin{equation}
        f(x)=\sum_{t\in T}x_tf(v_t).
    \end{equation}
    Ensuite le vecteur \( x - \sum_tx_tv_t\) est dans le noyau de \( f\), par conséquent nous le décomposons dans la base \( (u_s)\) :
    \begin{equation}
        x-\sum_tx_tv_t=\sum_{s\in S} x_su_s.
    \end{equation}
    Par conséquent
    \begin{equation}
        x=\sum_sx_su_s+\sum_tx_tv_t.
    \end{equation}

    En ce qui concerne la liberté nous écrivons
    \begin{equation}
        \sum_tx_tv_t+\sum_sx_su_s=0.
    \end{equation}
    En appliquant \( f\) nous trouvons que
    \begin{equation}
        \sum_tx_tf(v_t)=0
    \end{equation}
    et donc que les \( x_t\) doivent être nuls. Nous restons avec \( \sum_sx_su_s=0\) qui à son tour implique que \( x_s=0\).
\end{proof}
Un exemple d'utilisation de ce théorème en dimension infinie sera donné dans le cadre du théorème de Fréchet-Riesz, théorème~\ref{ThoQgTovL}.
\ifbool{isGiulietta}{Il existe une généralisation du théorème du rang pour les variétés différentiables en le théorème \ref{THOooSWKVooTJQsXc}.}{}

\begin{corollary}       \label{CORooCCXHooALmxKk}
    Soient deux espaces vectoriels \( E\) et \( F\) de même dimensions finies\footnote{Les deux mots sont importants : les dimensions doivent être égales et finies.}. Pour une application linéaire \( f\colon E\to F\), les trois conditions suivantes sont équivalentes :
    \begin{enumerate}
        \item
            \( f\) est injective;
        \item
            \( f\) est surjective;
        \item
            \( f\) est bijective.
    \end{enumerate}
\end{corollary}

\begin{proof}
    Si un endomorphisme \( f\colon E\to E\) est surjectif, alors \( \rang(f)=\dim(E)\), ce qui donne, par le théorème du rang~\ref{ThoGkkffA}, \( \dim\big( \ker(f) \big)=0\), c'est-à-dire que \( f\) est injectif.

    De la même façon, si \( f\) est injective, alors \( \dim\big( \ker(f) \big)=0\), ce qui donne \( \rang(f)=\dim(E)\) ou encore que \( f\) est surjective.
\end{proof}

\begin{example}
    Le corolaire \ref{CORooCCXHooALmxKk} n'est pas correct en dimension infinie. Par exemple en prenant \( f(e_1)=f(e_2)=e_1\) et ensuite \( f(e_k)=e_{k-1}\) pour tout \( k\geq 2\). Cette application est surjective mais pas injective.
\end{example}

Une conséquence du théorème du rang est que les endomorphismes ont un inverse à gauche et à droite égaux (lorsqu'ils existent). En résumé, ce que le corolaire \ref{CORooNFJLooJtzFwN} dit est que si \( f\circ g = \id\), alors \( g \circ f = \id\).
\begin{corollary}           \label{CORooNFJLooJtzFwN}
    Soit un endomorphisme \( f\) d'un espace vectoriel de dimension finie. Si \( f\) admet un inverse à gauche, alors
    \begin{enumerate}
        \item
            \( f\) est bijective,
        \item
            \( f\) admet également un inverse à droite,
        \item
            les inverses à gauche et à droite sont égaux.
    \end{enumerate}
    Tout cela tient également en remplaçant «gauche» par «droite».
\end{corollary}

\begin{proof}
    Soit \( g\), un inverse à gauche de \( f\) : \( gf=\id\). Cela implique que \( f\) est injective et que \( g\) est surjective, et donc qu'elles sont toutes deux bijectives par le corolaire~\ref{CORooCCXHooALmxKk}. Vu que \( f\) est bijective, elle admet également un inverse à droite, soit \( h\). Nous avons : \( gf=\id\) et \( fh=\id\).

    Alors \( gfh=h\) parce que \( gf=\id\), mais également \( gfh=g\) parce que \( fh=\id\). Donc \( g=h\).\footnote{C'est le même argument que celui employé pour la preuve du lemme~\ref{LEMooECDMooCkWxXf}~\ref{ITEMooOIWTooYqmMPP}, à ceci près que nous devions montrer l'existence de l'inverse à droite.}
\end{proof}
C'est ce corolaire qui nous permet d'écrire \( f^{-1}\) sans plus de précisions dès que \( f\) est une bijection.

\begin{example}[Pas en dimension infinie]
    Tout cela ne fonctionne pas en dimension infinie. Par exemple avec une base \( \{ e_k \}_{k\in \eN}\) nous pouvons considérer l'opérateur
    \begin{equation}
        f(e_k)=e_{k+1}.
    \end{equation}
    Il est injectif, mais pas surjectif. Si on pose
    \begin{equation}
        g(e_k)=\begin{cases}
            e_{k-1}    &   \text{si } k\geq 1\\
            0    &    \text{si } k=0
        \end{cases}
    \end{equation}
    alors nous avons \( gf=\id\), mais pas \( fg=\id\) parce que ce \( (fg)(e_0)=0\).
\end{example}

\begin{lemma}       \label{LEMooRZDTooEuLTrO}
    Si \( E\) et \( F\) sont des espaces vectoriels et si \( f\colon E\to F\) est une application linéaire inversible, alors son inverse est également linéaire.
\end{lemma}

\begin{proof}
    Nous avons \( f^{-1}(x+y)=f^{-1}(x)+f^{-1}(y)\). En effet,
    \begin{equation}
        f\big( f^{-1}(x)+f^{-1}(y) \big)=f\big( f^{-1}(x) \big)+f\big( f^{-1}(y) \big)=x+y.
    \end{equation}
    De la même façon,
    \begin{equation}
        f\big( \lambda f^{-1}(x) \big)=\lambda x, 
    \end{equation}
    donc \( f^{-1}(\lambda x)=\lambda f^{-1}(x)\).
\end{proof}

\begin{proposition}     \label{PROPooHLUYooNsDgbn}
    Soient un espace vectoriel \( E\) de dimension finie, un endomorphisme \( f\colon E\to E\) et une partie \( \{v_i\}_{i\in I}\) tel que \( \{f(v_i)\}_{i\in I}\) soit une base.

    Alors \( \{v_i\}_{i\in I}\) est une base.
\end{proposition}

\begin{proof}
    Soit \( x\in E\). Il existe une partie finie \( J\subset I\) et des scalaires \( \lambda_j\) tels que 
    \begin{equation}
        x=\sum_j\lambda_jf(v_j)=f\big( \sum_j\lambda_jv_j \big),
    \end{equation}
    ce qui prouve que \( f\) est surjective. Le corolaire \ref{CORooCCXHooALmxKk} nous dit alors que \( f\) est une bijection. L'application inverse est également linéaire par le lemme \ref{LEMooRZDTooEuLTrO}.

    Une application linéaire bijective (comme \( f^{-1}\)) transforme une base en une base par la proposition \ref{PROPooZFKZooBGLSex}. Donc 
    \begin{equation}
        f^{-1}\big( \{f(v_i)\} \big)
    \end{equation}
    est une base.
\end{proof}

\begin{proposition}     \label{PROPooADESooATJSrH}
    Soit un espace vectoriel \( E\) de dimension finie et deux applications linéaires \( f,g\colon E\to E\) telles que \( g\circ f=\id\). Alors \( f\) et \( g\) sont bijectives.
\end{proposition}

\begin{proof}
    En plusieurs étapes
    \begin{subproof}
        \item[\( f\) est injective]
            Si \( f(x)=f(y)\), alors en appliquant \( g\) nous avons 
            \begin{equation}
                g\big( f(x) \big)=g\big( f(y) \big),
            \end{equation}
            ce qui donne \( x=y\).
        \item[\( f\) est surjective]
            C'est maintenant le corolaire \ref{CORooCCXHooALmxKk}.
        \item[\( g\) est surjective]
            Pour tout \( x\in E\) nous avons \( g\big( f(x) \big)=x\). Donc l'image de \( f(E)\) par \( g\) est $E$. 
        \item[\( g\) est injective]
            C'est maintenant le corolaire \ref{CORooCCXHooALmxKk}.
    \end{subproof}
\end{proof}

%---------------------------------------------------------------------------------------------------------------------------
\subsection{Injection, surjection}
%---------------------------------------------------------------------------------------------------------------------------

\begin{definition}      \label{DEFooVTXWooVXfUnc}
    Soient deux espaces vectoriels \( E\) et \( V\). Une application \( f\colon E\to V\) est \defe{affine}{application affine} si il existe une application linéaire \( u\colon E \to V\) et un élément \( v\in V\) tel que
    \begin{equation}
        f(x)=u(x)+v
    \end{equation}
    pour tout \( x\in E\).
\end{definition}


\begin{definition}
    Soit un espace vectoriel \( V\). Une \defe{droite vectorielle}{droite vectorielle} dans \( V\) est un sous-espace vectoriel de dimension \( 1\). Un \defe{plan vectoriel}{plan vectoriel} est un sous-espace vectoriel de dimension \( 2\).

    Une partie \( D\) de \( V\) est une \defe{droite}{droite} si il existe \( v\in V\) tel que \( D-v\) soit une droite vectorielle.
\end{definition}

\begin{lemma}       \label{LEMooQQFFooEZYeck}
    Si \( D\) est une droite et si \( a,b\in D\), alors \( D-a=D-b\) et \( D-a\) est une droite vectorielle.
\end{lemma}

\begin{proof}
    Vu que \( D\) est une droite, il existe \( v\in V\) tel que \( D-v\) soit une droite vectorielle que nous notons \( L\). Nous allons montrer que \( D-a=D-v\). Vu que \( a\) est arbitraire, cela suffit.

    \begin{subproof}
        \item[\( D-a\subset D-v\)]
            Un élément de \( D-a\) est de la forme \( x-a\) avec \( x\in D\). Nous écrivons \( x-a\) sous la forme \( y-v\) et nous espérons que \( y\in D\). Allons-y : d'abord nous isolons \( y\) dans \( x-a=y-v\) :
            \begin{subequations}
                \begin{align}
                    y=x-a+v=(x-v)-(a-v)+v.
                \end{align}
            \end{subequations}
            Vu que \( x-v\) et \( a-v\) sont des éléments de \( L\), la somme est dans \( L\) et donc \( y=l+v\) pour un certain élément de \( l\in L\). Nous avons donc prouvé que \( y\in D\) et donc que \( x-a=y-v\in D-v\).
        \item[\( D-v\subset D-a\)]
            Nous notons \( x-v\) un élément générique que \( D-v\) (\( x\in D\)). En posant \( y-a=x-v\), nous trouvons
            \begin{equation}
                y=x-v+a=\underbrace{x-v}_{\in L}+\underbrace{(a-v)}_{\in L}+v
            \end{equation}
            Donc \( y\in D\) et \( x-v=y-a\in D-a\).
    \end{subproof}
\end{proof}

\begin{lemma}
    À propos de droites.
    \begin{enumerate}
        \item       \label{ITEMooYQCIooOrhRwj}
            Si \( L\) est une droite vectorielle, alors pour tout \( a\neq 0\) dans \( L\), nous avons \( L=\Image(f)\) où \( f\) est l'application linéaire donnée par
            \begin{equation}
                \begin{aligned}
                    f\colon \eK&\to V \\
                    \lambda&\mapsto \lambda a. 
                \end{aligned}
            \end{equation}
        \item       \label{ITEMooZIGMooGruFMP}
            Si \( D\) est une droite affine, alors pour tout \( a\neq b\) sur \( D\) nous avons \( D=\Image(f)\) où \( f\) est l'application affine donnée par
            \begin{equation}
                \begin{aligned}
                    g\colon \eK&\to V \\
                    \lambda&\mapsto a+\lambda(b-a). 
                \end{aligned}
            \end{equation}
    \end{enumerate}
\end{lemma}

\begin{proof}
    En deux parties.
    \begin{subproof}
        \item[Pour \ref{ITEMooYQCIooOrhRwj}]
            Vu que \( L\) est un sous-espace de dimension \( 1\), il possède une baser, disons \( \{ b \}\). En particulier \( a=\mu b\) pour un certain \( \mu\in \eK\). Si \( x\in L\) nous avons \( x=\lambda_x b\) pour un certain \( \lambda_x\), et donc
            \begin{equation}
                x=\frac{ \lambda_x }{ \mu }a.
            \end{equation}
            Donc \( x=f(\lambda_x/\mu)\). Cela prouve que \( L\subset\Image(f)\).

            L'inclusion inverse est simplement le fait que \( \lambda a\in L\) dès que \( a\in L\) parce que \( L\) est vectoriel.
        \item[Pour \ref{ITEMooZIGMooGruFMP}]
            Le lemme \ref{LEMooQQFFooEZYeck} nous indique qu'il existe une droite vectorielle \( L\) telle que \( D-x=L\) pour tout \( x\in D\).
            \begin{subproof}
                \item[\( D\subset\Image(g)\)]
                    Nous nommons \( f\colon \eK\to V\) l'application linéaire qui donne \( L\). Vu que \( b-a\in L\) nous avons
                    \begin{equation}
                        f(\lambda)=\lambda(b-a),
                    \end{equation}
                    et tout élément de \( L\) est de la forme \( f(\lambda)\). Nous avons aussi \( D=L+a\); donc un élément de \( D\) est de la forme \( f(\lambda)+a\) et donc de la forme \( \lambda(b-a)+a=g(\lambda)\).
                \item[\( \Image(g)\subset D\)]
                    Un élément de \( \Image(g)\) est de la forme \( a+\lambda(b-a)\) avec \( \lambda\in \eK\). Mais \( b-a\in L\), donc \( \lambda(b-a)\in L\) et 
                    \begin{equation}
                        g(\lambda)=a+\lambda(b-a)\in a+L=D.
                    \end{equation}
            \end{subproof}
    \end{subproof}
\end{proof}

\begin{example}
	Les exemples les plus courants d'applications affines sont les droites et les plans ne passant pas par l'origine.
	\begin{description}
		\item[Les droites] Une droite dans $\eR^2$ (ou $\eR^3$) qui ne passe pas par l'origine est l'image d'une fonction de la forme $s(t) =u t +v$, avec $t \in \eR$, et $u$ et $v$ dans $\eR^2$ ou $\eR^3$ selon le cas. 

		En choisissant des coordonnées adéquates, les droites peuvent être aussi vues comme graphes de fonctions affines. Dans le cas de $\eR^2$, on retrouve la fonction de l'exemple~\ref{ex_affine}, pour \( n = m = 1 \).

		\item[Les plans]
			De la même façon nous savons que tout plan qui ne passe pas par l'origine dans $\eR^3$ est le graphe d'une application affine, $P(x,y)= (a,b)^T\cdot(x,y)^T+(c,d)^T$, lorsque les coordonnées sont bien choisies.
	\end{description}
\end{example}

\begin{lemma}[\cite{ooEPEFooQiPESf}]        \label{LEMooDAACooElDsYb}
    Soit une application linéaire \( f\colon E\to F\).
    \begin{enumerate}
        \item       \label{ITEMooEZEWooZGoqsZ}
            L'application \( f\) est injective si et seulement s'il existe \( g\colon F\to E\) telle que \( g\circ f=\id|_E\).
        \item
            L'application \( f\) est surjective si et seulement s'il existe \( g\colon F\to E\) telle que \( f\circ g=\id|_F\).
    \end{enumerate}
\end{lemma}

\begin{proof}
    Nous démontrons séparément les deux affirmations.
    \begin{enumerate}
        \item
            Si \( f\) est injective, alors \( f\colon E\to \Image(f)\) est un isomorphisme. Si $V$ est un supplémentaire de \( \Image(f)\) dans \( F\) (c'est-à-dire \( F=\Image(f)\oplus V\)) alors nous pouvons poser \( g(x+v)=f^{-1}(x)\) où \( x+v\) est la décomposition (unique) d'un élément de \( F\) en \( x\in\Image(f)\) et \( v\in V\). Avec cela nous avons bien \( g\circ f=\id\).

            Inversement, s'il existe \( g\colon F\to E\) telle que \( g\circ f=\id\) alors \( f\colon E\to E\) doit être injective. Parce que si \( f(x)=0\) avec \( x\neq 0\) alors \( (g\circ f)(x)=0\neq x\).
        \item
            Si \( f\) est surjective nous pouvons choisir des éléments \( x_1,\ldots, x_p\) dans \( E\) tels que \( \{ f(x_i) \}\) soit une base de \( F\). Ensuite nous définissons
            \begin{equation}
                \begin{aligned}
                    g\colon F&\to E \\
                    \sum_k a_k f(x_k)&\mapsto \sum_k a_k x_k.
                \end{aligned}
            \end{equation}
            Cela donne \(  f\circ g=\id|_F\) parce que si \( v\in F\) alors \( v=\sum_kv_kf(x_k)\) avec \( v_k\in \eK\), et nous avons
            \begin{equation}
                (f\circ g)(v)=\sum_k v_k (f\circ g) \left(f(x_k)\right)
                             =f\left( \sum_k v_k x_k \right)
                             =\sum_k v_k f(x_k) = v.
            \end{equation}

            Inversement, s'il existe \( g\colon F\to E\) tel que \( f\circ g=\id\) alors \( f\) doit être surjective parce que
            \begin{equation}
                F=\Image(f\circ g)=f\big( \Image(g) \big)\subset \Image(f).
            \end{equation}
    \end{enumerate}
\end{proof}




%--------------------------------------------------------------------------------------------------------------------------- 
\subsection{Application linéaire associée à une matrice}
%---------------------------------------------------------------------------------------------------------------------------

Soient deux espaces vectoriels de dimension finie \( E,F\) sur le corps \( \eK\). Nous considérons les bases\footnote{C'est le théorème~\ref{ThonmnWKs} qui nous permet de considérer des bases. Et ce théorème ne fonctionne que parce que nous avons supposé une dimension finie.} \( \{ e_i \}\) pour \( E\) et \( \{ f_{\alpha} \}\) pour \( F\). 

\begin{definition}      \label{DEFooJVOAooUgGKme}
    Nous considérons l'application
    \begin{equation}        \label{EQooVZQWooMyFFeO}
        \begin{aligned}
            \psi\colon \eM(n\times m, \eK)&\to \aL(E,F) \\
            A&\mapsto f_A 
        \end{aligned}
    \end{equation}
    où \( f_A\) est définie par
    \begin{equation}        \label{EQooBVGHooJhFbMs}
        f_A(x)=\sum_{i\alpha}A_{\alpha i}x_if_{\alpha}
    \end{equation}
    si \( x_i\) sont les coordonnées de \( x\in E\) dans la base \( \{ e_i \}\).
\end{definition}



\begin{normaltext}
    Nous allons prouver un certain nombre de résultats montrant que cette application a toutes les propriétés imaginables permettant d'identifier les matrices aux applications linéaires : elle est un isomorphisme pour toutes les structures que vous pouvez raisonnablement imaginer.

    À cette application \( \psi\) il manque cependant une propriété importante : elle n'est pas canonique. Elle dépend des bases choisies. Autrement dit : nous avons à priori autant d'applications \( \psi\) différentes qu'il y a de choix de bases sur \( E\) et \( F\)\quext{Bonne question. Est-ce qu'il y a moyen de construire deux choix de bases donnant la même application \( \psi\) ? Écrivez-moi si vous savez la réponse.}. %TODOooFKSRooBtOZYc

    Nous allons prouver maintenant quelques résultats montrant que les matrices et les applications linéaires, dans le cas des espaces vectoriels \( \eK^n\) sont deux présentations de la même chose.
\end{normaltext}

\begin{normaltext}
    Lorsque \( A\in \eM(n,\eK)\) est une matrice et \( x\in \eK^n\) un vecteur, nous notons \( Ax\) l'élément de \( \eK^n\) donné par
    \begin{equation}        \label{EQooQFVTooMFfzol}
        (Ax)_i=\sum_jA_{ij}x_j.
    \end{equation}
    Autrement dit, \( Ax=f_A(x)\).

    Cette convention et de nombreuses autres à propos de matrice sera rappelée dans \ref{SECooBTTTooZZABWA}.
\end{normaltext}

\begin{propositionDef}      \label{PROPooGXDBooHfKRrv}
    Soient deux espaces vectoriels de dimension finie \( E,F\) sur le corps \( \eK\). Nous considérons les bases \( \{ e_i \}\) pour \( E\) et \( \{ f_{\alpha} \}\) pour \( F\). 

    Nous considérons l'application
    \begin{equation}
        \begin{aligned}
            \psi\colon \eM(n\times m, \eK)&\to \aL(E,F) \\
            A&\mapsto f_A 
        \end{aligned}
    \end{equation}
    où \( f_A\) est définie par
    \begin{equation}        \label{EQooZKEKooNYjvhP}
        f_A(x)=\sum_{i\alpha}A_{\alpha i}x_if_{\alpha}
    \end{equation}
    si \( x_i\) sont les coordonnées de \( x\in E\) dans la base \( \{ e_i \}\).

    Alors
    \begin{enumerate}
        \item       \label{ITEMooKZYYooZPTkpq}
            Nous avons
            \begin{equation}
                f_A(e_i)_{\alpha}=A_{\alpha i}.
            \end{equation}
        \item       \label{ITEMooANXFooGIuxUR}
            Nous avons
            \begin{equation}                \label{EQooOKOJooYgteNP}
                f_A(e_i)=\sum_{\alpha}A_{\alpha i}f_{\alpha}.
            \end{equation}
        \item       \label{ITEMooXLLLooKfigfB}
            Nous avons
            \begin{equation}        \label{EQooAXRJooUwHbjB}
                \big( f_A(x) \big)_{\alpha}=\sum_{i}A_{\alpha i}x_i.
            \end{equation}
        \item       \label{ITEMooHSMLooRJZref}
            L'application \( \psi\) est une bijection.
    \end{enumerate}
    Si \( f\) est une application linéaire, alors la matrice \( \psi^{-1}(f)\) est la \defe{matrice associée}{matrice d'une application linéaire} à \( f\) dans les bases choisies.
\end{propositionDef}

Remarque : les bases ne sont supposées être canoniques en aucun sens du terme. Les dimensions de \( E\) et \( F\) ne sont pas non plus supposées identiques.

\begin{proof}
    En nous rappelant que \( (e_j)_i=\delta_{ij}\) nous avons
    \begin{equation}        \label{EQooWGZHooIBoygB}
        f_A(e_j)=\sum_{i\alpha}A_{\alpha i}(e_j)_if_{\alpha}=\sum_{\alpha}A_{\alpha j}f_{\alpha},
    \end{equation}
    donc \( f_A(e_i)_{\alpha}=A_{\alpha i}\). Cela prouve la formule du point \ref{ITEMooKZYYooZPTkpq}.

    Le point \ref{ITEMooANXFooGIuxUR} est une simple somme sur \( \alpha\) de \ref{ITEMooKZYYooZPTkpq}.

    La formule \eqref{ITEMooXLLLooKfigfB} est simplement la composante \( f_{\alpha}\) de la définition \ref{EQooZKEKooNYjvhP}.

    Prouvons que \( \psi\) est injective. Si \( f_A=f_B\), nous avons en particulier \( f_A(e_i)_{\alpha}=f_B(e_i)_{\alpha}\) et donc \( A_{\alpha i}=B_{\alpha i}\).

    Prouvons que \( \psi\) est surjective. Pour cela nous considérons \( f\in \aL(E,F)\) et nous posons \( A_{\alpha i}=f(e_i)_{\alpha}\). Nous avons alors \( f=f_A\) parce que
    \begin{equation}
        f_A(x)=\sum_{i\alpha}A_{\alpha i}x_if_{\alpha}=\sum_{i\alpha}f(e_i)_{\alpha}x_if_{\alpha}=\sum_{\alpha}f(\sum_ix_ie_i)_{\alpha}f_{\alpha}=\sum_{\alpha}f(x)_{\alpha}f_{\alpha}=f(x).
    \end{equation}
\end{proof}

La proposition suivante montre que le produit matriciel correspond à la composition d'applications linéaires, pourvu que l'on travaille avec les bases canoniques sur \( \eK^n\).
\begin{proposition}[\cite{MonCerveau}]      \label{PROPooIYVQooOiuRhX}
    Soit un corps commutatif \( \eK\). Nous considérons des espaces vectoriels \( E\) et \( F\) munis de bases \( \{ e_i \}_{i=1,\ldots, n}\) et \( \{ f_{\alpha}\}_{\alpha\in 1,\ldots, m} \).

    L'application déjà définie\footnote{Notez la position du \( n\) et du \( m\). Sachez noter les bornes des sommes écrites dans la démonstration.}
    \begin{equation}
        \psi\colon \eM(m\times n,\eK)\to \aL(E,F)
    \end{equation}
    est un isomorphisme d'espaces vectoriels.
\end{proposition}

\begin{proof}
    Le fait que \( \psi\) soit une bijection est la proposition \ref{PROPooGXDBooHfKRrv}. Nous devons montrer que c'est linéaire. 

    Pour \( \lambda\in \eK\) nous avons le calcul
    \begin{equation}
        \psi(\lambda A)(e_k)=f_{\lambda A}(e_k)=\sum_{\alpha i}(\lambda A)_{\alpha i}\underbrace{(e_k)_i}_{=\delta_{ki}}f_{\alpha}=\lambda\sum_{\alpha}A_{\alpha k}f_{\alpha}=\lambda f_A(e_j).
    \end{equation}
    Donc \( \psi(\lambda A)=\lambda\psi(A)\).

    Si \( A,B\in \eM(n,\eK)\) nous avons de la même façon \( f_{A+B}=f_A+f_B\).
\end{proof}

\begin{proposition}     \label{PROPooCSJNooEqcmFm}
    Soient des espaces vectoriels \( E\), \( F\) et \( G\) de dimensions \( n\), \( m\) et \( p\) munis de bases\footnote{Avec trois, nous renonçons à utiliser des alphabets différents pour numéroter les éléments des bases.} \( \{ e_i \}\), \( \{ f_i \}\) et \( \{ g_i \}\). Nous considérons les deux applications
    \begin{equation}
        \psi\colon \eM(m\times n,\eK)\to \aL(E,F)
    \end{equation}
    et
    \begin{equation}
        \psi\colon \eM(p\times m,\eK)\to \aL(F,G).
    \end{equation}
    Nous avons 
    \begin{equation}
        f_A\circ f_B=f_{AB}
    \end{equation}
    pour toutes matrices \( A\in \eM(p\times m,\eK)\) et \( B\in \eM(m\times n,\eK)\).
\end{proposition}

\begin{proof}
    Nous considérons les applications linéaires associées à \( A\) et \( B\) : \( f_A\colon F\to G\) et \( f_B\colon E\to F\) et la composée \( f_A\circ f_B\colon E\to G\). Et puis c'est le calcul :
    \begin{subequations}
        \begin{align}
            (f_A\circ f_B)(e_k)&=f_A\big( \sum_{ij}B_{ij}(e_k)_jf_i \big)\\
            &=\sum_i B_{ik}f_A(f_i)\\
            &=\sum_iB_{ik}\sum_{rs}A_{rs}(f_i)_sg_r\\
            &=\sum_{ir}B_{ik}A_{ri}g_r\\
            &=\sum_r(AB)_{rk}g_r\\
            &=f_{AB}(e_k).
        \end{align}
    \end{subequations}
    Donc \( f_A\circ f_B=f_{AB}\) comme il se doit.
\end{proof}
    
Nous pouvons particulariser au cas où \( E=F=G\).
\begin{proposition}     \label{PROPooFMBFooEVCLKA}
    Si \( E\) est un espace vectoriel muni d'une base \( \{ e_i \}\), alors l'application
    \begin{equation}
        \psi\colon \eM(n,\eK)\to \End(E)
    \end{equation}
    est un isomorphisme d'algèbre\footnote{Définition \ref{DefAEbnJqI}.} et d'anneaux\footnote{Définition \ref{DEFooSPHPooCwjzuz}}.
\end{proposition}

\begin{proof}
    Le fait que \( \psi\) soit un isomorphisme d'algèbre est juste la combinaison entre les propositions \ref{PROPooIYVQooOiuRhX} et \ref{PROPooCSJNooEqcmFm}.

    En ce qui concerne l'isomorphisme d'anneaux, il faut en plus identifier les neutres. Le neutre pour la composition d'applications linéaires est l'application identité et le neutre pour la multiplication de matrices est la matrice identité. Nous devons donc montrer que \( \psi(\mtu)=\id\). Juste un calcul :
    \begin{equation}
        f_{\mtu}(x)=\sum_{ij}\delta_{ij}x_je_i=\sum_ix_ie_i=x.
    \end{equation}
    Donc oui, \( f_{\mtu}\) est l'application identité.
\end{proof}

Voilà. Soyez bien conscient que l'application \( \psi\) dont nous avons beaucoup parlé est surtout intéressante dans le cas des espaces de la forme \( \eK^n\). Dans ce cas, nous avons une identification canonique entre \( \eM(n,\eK)\) et \( \End(\eK^n)\) qui est un isomorphisme d'anneaux et d'algèbres.

Nous verrons que ce \( \psi\) respecte encore les inverses\footnote{Proposition \ref{PROPooNPMCooPmaCwu}.} et les déterminants\footnote{Proposition \ref{PROPooFKDXooKMSolt}.}.

\begin{normaltext}
    Il convient de ne pas confondre matrice et application linéaire (bien que nous le ferons sans vergogne). Une matrice est un bête tableau de nombres, tandis qu'une application linéaire est une application entre deux espaces vectoriels vérifiant certaines propriétés.

    Cependant si les espaces vectoriels \( E\) et \( F\) sont munis de bases, alors il y a une application
    \begin{equation}
        \psi\colon \eM(m\times n,\eK)\to \aL(E,F)
    \end{equation}
    qui a toutes les propriétés imaginables\footnote{Et elle en aura encore plus lorsque nous associerons les déterminants.}.

    Cette application dépend des bases choisies. Il n'y a donc pas de trucs comme «la matrice de telle application linéaire» ou comme «voici une matrice, nous considérons l'application linéaire associée». 

    Cependant, sur des espaces comme \( \eR^n\) ou plus généralement sur \( \eK^n\), nous avons une base canonique et toute personne raisonnable utilise toujours la base canonique (sauf mention du contraire). Dans ces cas il est sans danger de dire «la matrice associée à telle application linéaire» sans préciser les bases.

    Mais si un jour vous utilisez une base autre que la base canonique sur \( \eR^n\), précisez-le et plutôt deux fois qu'une\footnote{Au passage, non, les coordonnées polaires ne sont pas une base de \( \eR^2\). C'est un système de coordonnées, et ce n'est pas la même chose.}.
\end{normaltext}

\begin{normaltext}
    Tant que nous sommes à parler de matrice et d'applications linéaires, les plus acharnés anti-abus de language\footnote{Dont l'auteur de ces lignes fait partie.} remarqueront qu'il n'est pas vrai que «étant donné une base, une application linéaire a une matrice».

    En effet, une base est une partie libre et génératrice (définition \ref{DEFooNGDSooEDAwTh}). Or une partie d'un ensemble n'est pas muni d'un ordre. Toutes les permutations de colonnes de la matrice sont encore possibles d'après l'ordre que l'on met sur les vecteurs de la base.

    Encore une fois, la base canonique n'a pas de problèmes parce que les \( \{ e_i \}\) de \( \eR^n\) viennent avec un ordre indiscutable. Plus généralement, très souvent, lorsqu'on construit une base, la construction suggère un ordre.
\end{normaltext}


\begin{proposition}     \label{PROPooNPMCooPmaCwu}
    Une matrice est inversible si et seulement si son application linéaire associée est inversible. Dans ce cas, nous avons
    \begin{equation}
        f_A^{-1}=f_{A^{-1}}.
    \end{equation}
\end{proposition}

\begin{proof}
    Dans le sens direct, si \( A\) est inversible nous avons \( AA^{-1}=\mtu\). Donc
    \begin{equation}        \label{EQooQQOSooBKVqXh}
        f_A\circ f_{A^{-1}}=f_{AA^{-1}}=f_{\mtu}=\id
    \end{equation}
    où nous avons utilisé la proposition \ref{PROPooCSJNooEqcmFm} pour la composition et la proposition \ref{PROPooFMBFooEVCLKA} pour l'identité. L'égalité \eqref{EQooQQOSooBKVqXh} indique que \( f_A\) est inversible et que son inverse est \( f_{A^{-1}}\).

    Dans l'autre sens, l'application \( f_A^{-1}\) existe. Soit \( B\in \eM(n,\eK)\) sa matrice. Alors nous avons
    \begin{equation}
        f_{\delta}=\id=f_A\circ f_B=f_{AB}.
    \end{equation}
    Le fait que l'application \(\psi\colon A\to f_A\) soit une bijection\footnote{Proposition \ref{PROPooGXDBooHfKRrv}\ref{ITEMooHSMLooRJZref}.} implique que \( AB=\mtu\), c'est-à-dire que\footnote{Corollaire~\ref{CORooBQLXooTeVfgb}.} \( A\) est inversible et que \( B=A^{-1}\).
\end{proof}


Parlons de rangs de matrices à présent. Nous donnons une définition qui est clairement inspirée de l'application linéaire associée.
\begin{lemma}
  Soit \( A \) une matrice de rang \( r \), et \( G \) une matrice de tranformation telle que définie en~\ref{DEF_transvectionDilatationPermutation}. Le rang de la matrice \( GA \) est alors aussi \( r \).
\end{lemma}

\begin{proof}
Il suffit de collecter les renseignements suivants: proposition~\ref{PROPooXUFKooOaPnna} point~\ref{ITEMooRangsDilatations} pour les dilatations; proposition~\ref{PROPooFQRDooRPfuxk} point~\ref{ITEMooRangsPermutations} pour les permutations; proposition~\ref{PROPooPYNHooLbeVhj} point~\ref{ITEMooRangsTransvections} pour les transvections.
\end{proof}

\begin{proposition}[\cite{ooKBOMooSkKHvu}]       \label{DEFooCSGXooFRzLRj}
    Le rang d'une matrice \(A \in \eM(n,\eK)\) est la dimension de la partie de \( \eK^n\) engendrée par ses colonnes.
\end{proposition}
\index{rang!d'une matrice}
\begin{proof}
  On note \( v_1,\dots, v_n\) les vecteurs qui forment les colonnes de \( A \), et \( d \) la dimension de l'espace qu'ils engendrent. Quitte à renuméroter les vecteurs\footnote{Formellement, à coup de permutation de colonnes de \(A \), ce qui ne change pas le rang.}, on peut dire que tous les \( v_k \), avec \( d < k \leq n \), sont combinaisons linéaires des \( v_i \) pour \( i < d\). Il s'ensuit qu'un mineur de taille \( d+1 \) (ou plus grand) dans \( A \) est nécessairement nul, et qu'un mineur de taille \( d \) constitué de composantes des \( v_i \) est nécessairement non nul.
\end{proof}


Lorsque \( A\) est une matrice, nous notons \( f_A\) l'application linéaire associée à la matrice \( A\) par l'application \eqref{EQooVZQWooMyFFeO}.

\begin{lemma} \label{LEMVecsaRgFixe}
    Soit \( \eK \) un corps commutatif\footnote{Comme toujours.}. Si \( A \) est une matrice carrée d'ordre \( n \) et de rang \( r \) à coefficients dans \( \eK \), alors il existe des vecteurs \( (x_i)_{i=1,\dots,n} \) formant une base de \( \eK^n \) tels que 
    \begin{equation}
        f_A(x_i)\neq 0
    \end{equation}
    pour \( i\leq r\) et
    \begin{equation}
        f_A(x_i) = 0
    \end{equation}
    pour \( i > r \).
\end{lemma}

\begin{proof}
    Soit \( V\) le sous-espace de \( \eK^n\) engendré par les colonnes de \( A\). Nous considérons la base canonique \( \{ e_i \}\) de \( \eK^n\), ainsi que \( v_i\) le vecteur créé par la \( i\)\ieme colonne de \( A\). Nous avons
    \begin{equation}
        v_i=f_A(e_i).
    \end{equation}
    Les vecteurs \( v_i\) engendrent \( V\), donc nous pouvons en extraire une base par le théorème \ref{ThoMGQZooIgrXjy}\ref{ITEMooTZUDooFEgymQ}. Soit donc \( \{ v_j \}_{i\in J}\) une base de \( V\) avec \( J\subset\{ 1,\ldots, n \}\).

    La base de \( \eK^n\) que nous cherchons commence par les vecteurs \( \{ e_j \}_{j\in J}\). Ces vecteurs vérifient \( f_A(e_j)=v_j\neq 0\) parce que des vecteurs d'une base ne sont jamais nuls.

    % Note : toute la ligne suivante fait des références qui peuvent être vers le futur parce que ce sont des choses qui ne sont 
    %        pas utilisées dans la démonstration.
    Pour la suite de la base, nous pourrions penser au théorème de la base incomplète\footnote{Théorème \ref{ThonmnWKs}\ref{ITEMooFVJXooGzzpOu}.}, mais les vecteurs ainsi complétant la base ne sont pas garantis de s'annuler par \( f_A\). Voir l'exemple \ref{EXooRKVQooZOGDEf}.

    L'idée est d'utiliser le noyau de \( f_A\) qui est un sous-espace vectoriel par la proposition \ref{PROPooRLLPooKYzsJp}. Soit une base\footnote{Cette base contient \( n-r\) éléments, mais ce n'est pas très important pour la suite.} \(  \{ z_k \}  \) de \( \ker(f)\). Les vecteurs \( \{ e_j \}_{j\in J}\) forment une base de \( \Image(f_A)\). Vu que les \( z_i\) forment une base de \( \ker(f_A)\), le théorème du rang \ref{ThoGkkffA} dit alors que \( \{ e_j \}_{j\in J}\cup \{ z_k \}\) est une base de \( \eK^n\).

    Il y a \( r\) éléments dans \( J\) parce que l'espace engendré par les colonnes de \( A\) est de dimension \( r\) par hypothèse. Donc il y a \( n-r\) éléments dans les \( z_k\) pour que le tout ait le bon nombre d'éléments.
\end{proof}

\begin{example}     \label{EXooRKVQooZOGDEf}
    Soit la matrice 
    \begin{equation}
        A=\begin{pmatrix}
            1    &   1    \\ 
            2    &   2    
        \end{pmatrix}.
    \end{equation}
    Elle est de rang \( 1\). En suivant l'idée de la démonstration, nous commençons la base de \( \eR^2\) par le vecteur \( e_1\) qui vérifie
    \begin{equation}
        f_A(e_1)=\begin{pmatrix}
            1    \\ 
            2    
        \end{pmatrix}.
    \end{equation}
    L'utilisation du théorème de la base incomplète ne permet pas de trouver un second vecteur de base \( v\) tel que \( f_A(v)=0\). En effet ce théorème donne juste l'existence d'une complétion de la base, mais pas de propriétés particulières de la base obtenue. Elle pourrait donner \( v=e_2\) comme second vecteur de base. Mais alors
    \begin{equation}
        f_A(v)=f_A(e_2)=\begin{pmatrix}
            1    \\ 
            2    
        \end{pmatrix}\neq 0.
    \end{equation}

    Au contraire, le noyau de \( f_A\) est donné par le sous-espace engendré par \( \begin{pmatrix}
        1    \\ 
        -1    
    \end{pmatrix}\). Une base convenable est donc \( \{ e_1, e_1-e_2 \}\).
\end{example}


\begin{proposition} \label{PropEJBZooTNFPRj}
    Si \( A\) est la matrice d'une application linéaire, alors le rang de cette application linéaire est égal au rang de \( A \), c'est-à-dire à la taille de la plus grande matrice carrée de déterminant non nul contenue dans \( A\).
\end{proposition}

\begin{proof}
  Notons \( E \) et \( F \) les espaces vectoriels concernés, de dimensions respectives \( n \) et \( m \) et de bases fixées \( (e_i)_{i=1}^n \) et \( (f_j)_{j=1}^m \). Soit \( T: E \to F \) une application linéaire. Ce choix fait, la proposition nous demande d'avoir \( A \)  associée à \( T \), en d'autres termes: \( (T(e_i))_j = A_{ij} \). Donc en fait les \( T(e_i) \) sont les colonnes de \( A \). Le rang de \( T \) est la dimension de l'espace vectoriel engendré par les \( T(e_i) \), donc la dimension de l'espace vectoriel engendré par les colonnes de \( A \), donc en fait le rang de \( A \) par la proposition~\ref{DEFooCSGXooFRzLRj}.
\end{proof}


\begin{proposition}     \label{PROPooEGNBooIffJXc}
    Le rang d'une application linéaire est égal au rang de sa matrice dans n'importe quelle base.
\end{proposition}

\begin{proof}
  Soit \( T: E \to F \) une application linéaire. Quelles que soient les bases choisies, l'espace vectoriel \( W = \Im(T) \) a toujours la même dimension. Les vecteurs qui constituent les colonnes des matrices sont des éléments de \( W \) qui le génèrent.
\end{proof}

---------------------------------------------------------------------------------------------------------------------------
\subsection{Matrices équivalentes et semblables}
%---------------------------------------------------------------------------------------------------------------------------

\begin{definition}  \label{DefBLELooTvlHoB}
    Deux relations d'équivalence entre les matrices.
    \begin{enumerate}
        \item   \label{ItemPFXCooOUbSCt}
    Deux matrices \( A\) et \( B\) sont \defe{équivalentes}{matrice!équivalence} dans \( \eM(n,\eK)\) s'il existe \( P,Q\in\GL(n,\eK)\) telles que \( A=PBQ^{-1}\).
\item
    Deux matrices sont \defe{semblables}{matrices!similitude} s'il existe une matrice \( P\in \GL(n,\eK)\) telle que \( A=PBP^{-1}\).
    \end{enumerate}
\end{definition}

\begin{lemma}   \label{LemZMxxnfM}
    Une matrice de rang\footnote{Définition~\ref{DEFooVVBYooJbliTi}.} \( r\) dans \( \eM(n,\eK)\) est équivalente à la matrice par blocs
    \begin{equation}
        J_r=\begin{pmatrix}
            \mtu_r    &   0    \\
            0    &   0
        \end{pmatrix}.
    \end{equation}
\end{lemma}
\index{rang!classe d'équivalence}

\begin{proof}
    Nous devons prouver que pour toute matrice \( A\in\eM(n,\eK)\) de rang \( r\), il existe \( P,Q\in\GL(n,\eK)\) telles que \(QAP=J_r\). Soit \( \{ e_i \}\) la base canonique de \( \eK^n\), puis \( \{ f_i \}\) une base telle que \( Af_i=0\) dès que \( i>r\), qui existe par le lemme~\ref{LEMVecsaRgFixe}.

    Nous considérons la matrice inversible \( P\) telle que \( Pe_i=f_i\); ses colonnes sont donc précisément les \( f_i \), si bien que
    \begin{equation}
        APe_i=Af_i=\begin{cases}
            0    &   \text{si } i>r\\
            \neq 0    &    \text{sinon}.
        \end{cases}
    \end{equation}
    La matrice \( AP\) se présente donc sous la forme
    \begin{equation}
        AP=\begin{pmatrix}
            M    &   0    \\
            *    &   0
        \end{pmatrix}
    \end{equation}
    où \( M\) est une matrice \( r\times r\). Nous considérons maintenant une base \( \{ g_i \}_{i=1,\ldots, n}\) dont les \( r\) premiers éléments sont les \( r\) premières colonnes de \( AP\) et une matrice inversible \( Q\) telle que \( Qg_i=e_i\). Alors
    \begin{equation}
        QAPe_i=\begin{cases}
            e_i    &   \text{si } i<r\\
            0    &    \text{sinon}.
        \end{cases}.
    \end{equation}
    Cela signifie que \( QAP\) est la matrice \( J_r\).
\end{proof}

\begin{corollary}[Équivalence et rang]      \label{CorGOUYooErfOIe}
    Deux matrices sont équivalentes\footnote{Définition~\ref{DefBLELooTvlHoB}\ref{ItemPFXCooOUbSCt}.} si et seulement si elles sont de même rang.
\end{corollary}

\begin{proof}
    D'abord il y a des implicites dans l'énoncé. Vu que nous voulons soit par hypothèse soit par conclusion que les matrices \( A\) et \( B\) soient équivalentes, nous supposons qu'elles ont même dimension. Soient donc \( A\) et \( B\) deux matrices carrées d'ordre \( n \).

    Par le lemme~\ref{LemZMxxnfM}, deux matrices de même rang \( r\) sont équivalentes à \( J_r\). Elles sont donc équivalentes entre elles.

    Inversement, supposons que \( A\) et \( B\) soient deux matrices équivalentes : \( A=PBQ^{-1}\) avec \( P\) et \( Q\) inversibles. Alors
    \begin{subequations}
        \begin{align}
            \Image(PBQ^{-1})&=\{ PBQ^{-1}v\tq v\in \eK^n \}\\
            &=PB\underbrace{\{ Q^{-1}v\tq v\in \eK^n \}}_{=\eK^n}\\
            &=P\big( B(\eK^n) \big).
        \end{align}
    \end{subequations}
    L'ensemble \( B(\eK^n)\) est un sous-espace vectoriel de \( \eK^n\). Vu que le rang de \( P\) est maximum, la dimension de \( P\big( B(\eK^n) \big)\) est la même que celle de \( B(\eK^n)\). Par conséquent
    \begin{equation}
        \dim\Big( \Image(PBQ^{-1}) \Big)=\dim\big( B(\eK^n) \big)=\rang(B).
    \end{equation}
    Le membre de gauche de cela n'est autre que \( \rang(A)=\dim\big( \Image(PBQ^{-1}) \big)\).
\end{proof}
